\section*{Hoofdstuk \ref{ch:b1:organomagnesium} --- Organomagnesiumreagentia}

\begin{solution}{exo:b1:organomagnesium:1}
\ce{CH3CH2Br + Mg -> CH3CH2MgBr}; \ce{C6H5Br + Mg -> C6H5MgBr} (in droge
ether). C--Br: Br (2.96) is elektronegatiever dan C (2.55), koolstof is
positief. C--Mg: Mg (1.31) is minder elektronegatief, koolstof is
negatief.
\end{solution}

\begin{solution}{exo:b1:organomagnesium:2}
Water: \ce{CH3MgI + H2O -> CH4 + Mg(OH)I}.
Ethanol: \ce{CH3MgI + CH3CH2OH -> CH4 + CH3CH2OMgI}.
Ethaanzuur: \ce{CH3MgI + CH3COOH -> CH4 + CH3COOMgI}.
Deuteriumoxide: \ce{CH3MgI + D2O -> CH3D + Mg(OD)I}.
\end{solution}

\begin{solution}{exo:b1:organomagnesium:3}
Methanal: propaan-1-ol (primair). Ethanal: butaan-2-ol (secundair).
Propanon: 2-methylbutaan-2-ol (tertiair). Benzaldehyde:
1-fenylpropaan-1-ol (secundair). Koolstofdioxide: propaanzuur.
\end{solution}

\begin{solution}{exo:b1:organomagnesium:4}
Zijn OH-groep is zuur genoeg om het reagens te vernietigen zodra het
ontstaat: het molecuul zou zichzelf (of zijn buren) protoneren; zo
ontstaat het magnesium\-alkoxide van butaan-1-ol, \ce{CH3(CH2)3OMgBr}, en
dat is geen bruikbaar reagens. De OH moet eerst beschermd worden
(\cref{ch:b1:acetals-protection} toont het idee voor een aldehyde).
\end{solution}

\begin{solution}{exo:b1:organomagnesium:5}
Het C--Mg-paar valt het carbonylkoolstofatoom van ethanal aan, waarbij het
$\pi$-paar van C=O naar O gaat: \ce{CH3CH(CH3)O^-} met \ce{MgBr+}; daarna
geeft \ce{H3O+} een proton aan de zuurstof: propaan-2-ol. Propaan-2-ol
draagt twee methylgroepen op het carbinolkoolstofatoom: het is \omterm{def:b1:stereochemistry-in-depth:stereocentre}{achiraal}.
(Met een andere R, bijvoorbeeld ethylmagnesiumbromide op ethanal, is het
product butaan-2-ol \omterm{def:b1:stereochemistry-in-depth:stereocentre}{chiraal} maar racemisch: het vlakke carbonyl wordt even
gemakkelijk aan beide zijden aangevallen.)
\end{solution}

\begin{solution}{exo:b1:organomagnesium:6}
Het carbinolkoolstofatoom draagt één methyl- en twee ethylgroepen:
ethylmagnesiumbromide met butaan-2-on, of methylmagnesiumbromide met
pentaan-3-on.
\end{solution}

\begin{solution}{exo:b1:organomagnesium:7}
Bereid \ce{(CH3)2CHMgBr} uit 2-broompropaan en magnesium in droge ether,
giet het op vast \ce{CO2}, en zuur daarna aan: \ce{(CH3)2CHCOOH}.
$n = 12.3/123.0 = 0.100$ mol;
$0.100 \times 0.70 \times 88.0 = \qty{6.2}{g}$.
\end{solution}

\begin{solution}{exo:b1:organomagnesium:8}
$0.18/18.0 = \qty{0.010}{mol}$ water vernietigt \qty{0.010}{mol}
reagens: \qty{20}{\percent}; er ontstaat benzeen,
\ce{C6H5MgBr + H2O -> C6H6 +
Mg(OH)Br}.
\end{solution}

\begin{solution}{exo:b1:organomagnesium:9}
Het reagens, een \omterm{def:b1:electronic-effects:nucleophile}{nucleofiel}, reageert met broombenzeen dat nog aanwezig
is: in totaal \ce{C6H5MgBr + C6H5Br -> C6H5-C6H5 + MgBr2}. Het bromide
langzaam toevoegen aan een overmaat magnesium houdt zijn concentratie
laag: het ontmoet eerder vers metaal dan reagens dat al gevormd is.
\end{solution}

\begin{solution}{exo:b1:organomagnesium:10}
1-Fenylethaan-1-ol: \ce{CH3MgBr} (uit broommethaan) met benzaldehyde, of
\ce{C6H5MgBr} (uit broombenzeen) met ethanal. 2-Fenylpropaan-2-ol:
\ce{C6H5MgBr} met propanon, of \ce{CH3MgBr} met
1-fenylethaan-1-on.
\end{solution}

\begin{solution}{exo:b1:organomagnesium:11}
Het keton wordt aan het reagens toegevoegd, zodat het reagens, bereid in
de droge kolf, die nooit verlaat en de warmte die elke portie vrijmaakt
door de kokende ether opgenomen wordt. Alleen water zou het geleiachtige
basische zout \ce{Mg(OH)Br} geven, dat het product insluit; koud verdund
zuur lost het op als \ce{Mg^2+}, en de koude beperkt nevenreacties van de
tertiaire alcohol in zuur milieu.
\end{solution}

\begin{solution}{exo:b1:organomagnesium:12}
Elke mol reagens geeft met water één mol basisch magnesiumzout,
\ce{RMgBr + H2O -> RH + Mg(OH)Br}, waarvan de \ce{OH-} door het zuur
getitreerd wordt: de hoeveelheid reagens is gelijk aan het verbruikte
zuur, \qty{0.088}{mol}, een \omterm{def:b1:extent-q-and-k:fractional-extent}{rendement} van \qty{88}{\percent}.
\end{solution}

\begin{solution}{pb:b1:organomagnesium:1}
\textbf{1.} \ce{C6H5Br + Mg -> C6H5MgBr}.
\textbf{2.} Mg $0.535/24.3 = \qty{0.0220}{mol}$; \ce{C6H5Br} $3.14/156.9 =
\qty{0.0200}{mol}$: magnesium in lichte overmaat.
\textbf{3.} Water vernietigt het reagens: \ce{C6H5MgBr + H2O -> C6H6 +
Mg(OH)Br}.
\textbf{4.} Ze houdt de waterdamp van de lucht tegen, die anders via de
koeler binnenkomt.
\textbf{5.} De ether geeft \omterm{def:b1:lewis-resonance-vsepr:lewis}{vrije paren} aan magnesium (\omterm{def:b1:lewis-resonance-vsepr:lewis-acid}{lewisbase}): zonder
ether ontstaat het reagens niet of blijft het niet in oplossing.
\textbf{6.} \ce{C6H5MgBr + C6H5Br -> C6H5-C6H5 + MgBr2}; bifenyl
$0.15/154.0 = \qty{9.7e-4}{mol}$, dat \qty{9.7e-4}{mol} broombenzeen en
evenveel reagens verbruikt heeft: in totaal \qty{1.9e-3}{mol} broombenzeen
verloren.
\textbf{7.} $3.28/182.0 = \qty{0.0180}{mol}$.
\textbf{8.} Het C--Mg-paar valt het carbonylkoolstofatoom aan, het
$\pi$-paar gaat naar zuurstof: \ce{(C6H5)3C-O^-} \ce{MgBr+}.
\textbf{9.} Hoogstens beschikbaar reagens $0.0200 - 0.0019 = \qty{0.0181}{mol}$;
benzofenon \qty{0.0180}{mol}: benzofenon is beperkend, net.
\textbf{10.} Nee: het carbinolkoolstofatoom draagt drie identieke
fenylgroepen.
\textbf{11.} Het reagens blijft in zijn droge kolf, en de warmte van elke
portie keton wordt door de kokende ether opgenomen.
\textbf{12.} Het zou een deel van het reagens vernietigen (benzeen) en
benzofenon onomgezet laten.
\textbf{13.} \ce{(C6H5)3COMgBr + H3O+ -> (C6H5)3COH + Mg^2+ + Br- + H2O}.
\textbf{14.} Het zuur lost de basische magnesiumzouten op die alleen
water zou laten neerslaan; de koude beperkt de vorming van het \omterm{def:b1:electronic-effects:intermediates}{carbokation}
van vraag 18.
\textbf{15.} Trifenylmethanol en bifenyl in de etherlaag; de
magnesiumzouten in de waterlaag.
\textbf{16.} De vaste stof wassen (of fijnwrijven) met koude
petroleumether: bifenyl lost op, trifenylmethanol blijft achter.
\textbf{17.} Een scherp smeltpunt gelijk aan het getabelleerde (of één
enkele vlek bij dunnelaagchromatografie).
\textbf{18.} \ce{(C6H5)3COH + H+ -> (C6H5)3C+ + H2O}: een tertiair
\omterm{def:b1:electronic-effects:intermediates}{carbokation} waarvan de lading door resonantie over drie benzeenringen
gespreid is, zeer stabiel, en gekleurd door die uitgebreide conjugatie.
\textbf{19.} $3.75/260.0 = \qty{0.0144}{mol}$.
\textbf{20.} $0.0144/0.0180 = \qty{80}{\percent}$.
\textbf{21.} Verliezen bij het overgieten en de herkristallisatie;
reagens vernietigd door sporen vocht; onvolledige reactie.
\textbf{22.} Benzoëzuur, \ce{C6H5COOH}, na aanzuren.
\textbf{23.} $0.0181 \times 122.0 = \qty{2.21}{g}$.
\textbf{24.} \ce{C6H5MgBr + CO2 -> C6H5COOMgBr}, daarna
\ce{C6H5COOMgBr + H3O+ -> C6H5COOH + Mg^2+ + Br- + H2O}.
\textbf{25.} \textbf{\qty{80}{\percent}}.
\end{solution}
